% Original applied explanation; see Chapter04/SOURCES.md.
\subsubsection{মোবাইলে দোকানের পেজ ব্যবহার করা}
কম্পিউটারে দোকানের পণ্যের ছবি ও বিবরণ পাশাপাশি রাখা যায়। মোবাইলে জায়গা কম হলে সেগুলো ওপর-নিচে রাখা সুবিধাজনক। Responsive বিন্যাস পর্দার জায়গা অনুযায়ী এই সাজানো বদলায়; ছবি ছোট করাই তার একমাত্র কাজ নয়।

দর্শকের প্রথম কাজ ঠিক করবে: পণ্যের নাম, দাম ও যোগাযোগ যেন খুঁজে পাওয়া যায়। ওয়্যারফ্রেমে সেই তথ্যের স্থান চিহ্নিত করবে, HTML-এ অর্থপূর্ণ গঠন লিখবে এবং CSS দিয়ে মাপ ও বিন্যাস ঠিক করবে। ছোট পর্দায় পড়া, লিংক চাপা ও তথ্যের ক্রম পরীক্ষা করলে পরিকল্পনার সমস্যা ধরা পড়ে।
